%% toc: P15 171 171.1,171.2,171.3,171.4
\section{Statistical learning: proofs}
\label{sec:annex-p15-w02b-p15-064-prediction-risk}

Throughout, the prediction problem, the risk \(R\), the empirical risk
\(\widehat R_n\), \(R_{\mathcal F}\) and \(R^{*}\) are those of
\cref{def:p15-w02b-p15-064-prediction-problem,def:p15-w02b-p15-064-risk}.
For a random element \(W\) with values in a measurable space, \(C_W\) denotes
a selected version of conditional expectation given \(\sigma(W)\), as in
\cref{def:p15-general-conditional-expectation}.

\begin{annexitem}{annex:p15-w02b-p15-064-empirical-estimate}{thm:p15-w02b-p15-064-empirical-estimate}{Unbiasedness, the second-moment bound and the finite-class consequence}
\begin{proof}
\paragraph{Unbiasedness.}
Put \(L_i=\ell(X_i,Y_i,f(X_i))\) and \(L=\ell(X,Y,f(X))\). Since \(f\) does
not depend on the sample, \(L_i\) is the composition of the record
\((X_i,Y_i)\) with the fixed measurable map
\((x,y)\mapsto\ell(x,y,f(x))\); it is nonnegative and has the law of \(L\),
so \(EL_i=EL=R(f)\in[0,\infty]\). Linearity of the integral for nonnegative
extended variables gives
\(E\,\widehat R_n(f)=n^{-1}\sum_iEL_i=R(f)\).

\paragraph{One predictor.}
Assume \(\sigma^2(f)<\infty\), that is \(EL^2<\infty\); then
\(R(f)=EL<\infty\). Put \(D_i=L_i-R(f)\) and
\(D=\widehat R_n(f)-R(f)=n^{-1}\sum_iD_i\). The \(D_i\) are measurable
functions of the independent records, hence independent; each has mean zero
and \(ED_i^2=\sigma^2(f)\). For \(i\neq j\), \(D_i\) and \(D_j\) are
independent and integrable, so \(E(D_iD_j)=ED_i\,ED_j=0\). Expanding the
square,
\[
 ED^2=\frac1{n^2}\sum_{i=1}^nED_i^2=\frac{\sigma^2(f)}n.
\]
Pointwise \(\mathbf1_{\{|D|\geq\varepsilon\}}\leq D^2/\varepsilon^2\);
taking expectations gives the displayed bound.

\paragraph{Finite class.}
Let \(D^{(j)}=\widehat R_n(f_j)-R(f_j)\). Each \(f_j\) is fixed before
sampling, so the previous paragraph applies to each. The event
\(\{\max_j|D^{(j)}|\geq\varepsilon\}\) is the union of the \(r\) events
\(\{|D^{(j)}|\geq\varepsilon\}\); its probability is at most the sum of
theirs, which is the displayed bound.

\paragraph{The selected predictor.}
On the complementary event, \(|D^{(j)}|<\varepsilon\) for every \(j\). Let
\(f^{\circ}\in\mathcal F\) attain \(R_{\mathcal F}=\min_jR(f_j)\); a minimum
of finitely many finite numbers is attained. A \(\gamma\)-approximate
minimizer \(\widehat f\in\mathcal F\) satisfies
\(\widehat R_n(\widehat f)\leq\min_j\widehat R_n(f_j)+\gamma
\leq\widehat R_n(f^{\circ})+\gamma\). Since both \(\widehat f\) and
\(f^{\circ}\) are members of \(\mathcal F\), their deviations are below
\(\varepsilon\), whatever member \(\widehat f\) is. Therefore
\[
 R(\widehat f)<\widehat R_n(\widehat f)+\varepsilon
 \leq\widehat R_n(f^{\circ})+\gamma+\varepsilon
 <R(f^{\circ})+2\varepsilon+\gamma
 =R_{\mathcal F}+2\varepsilon+\gamma.
\]
\end{proof}
\end{annexitem}

\begin{annexitem}{annex:p15-w02b-p15-064-bounded-rescaling}{thm:p15-w02b-p15-064-bounded-rescaling}{Rescaling a bounded loss}
\begin{proof}
Put \(\ell'(h,z)=\ell(h,z)/B\). It is measurable with values in \([0,1]\);
the class, the sample and its independence are unchanged, so \(\ell'\)
satisfies \cref{def:p15-conc-learning}. Its risks and empirical risks are
\(R'=R/B\) and \(\widehat R'_n=\widehat R_n/B\). By
\cref{thm:p15-conc-finite-learning}, there is an event of probability at
least \(1-\delta\) on which \(|\widehat R'_n(h)-R'(h)|\leq b_n\) for every
\(h\), and on which every \(\eta'\)-approximate minimizer for \(\ell'\)
satisfies \(R'(\widehat h)\leq\min_hR'(h)+2b_n+\eta'\); the event is the
uniform deviation event and does not depend on the tolerance \(\eta'\).
Multiplying the first inequality by \(B\) gives
\(|\widehat R_n(h)-R(h)|\leq Bb_n\). An \(\eta\)-approximate minimizer for
\(\ell\), meaning \(\widehat R_n(\widehat h)\leq\min_h\widehat R_n(h)+\eta\),
is after division by \(B\) an \(\eta/B\)-approximate minimizer for \(\ell'\).
Taking \(\eta'=\eta/B\) and multiplying by \(B\) gives
\(R(\widehat h)\leq\min_hR(h)+2Bb_n+\eta\) on the same event.
\end{proof}
\end{annexitem}

\begin{annexitem}{annex:p15-w02b-p15-064-squared-bayes}{thm:p15-w02b-p15-064-squared-bayes}{Orthogonality of the residual to every function of the input}
\begin{proof}
\paragraph{The regression function.}
Since \(EY^2<\infty\), \(Y\) is integrable, and
\cref{thm:p15-conditional-factorization} gives a measurable
\(m:\mathcal X\to\mathbb R\) with \(m(X)\) a version of \(C_XY\). By
\cref{thm:p15-conditional-lp} with \(p=2\),
\(E\,m(X)^2\leq EY^2<\infty\). Hence \(U=Y-m(X)\) satisfies
\(EU^2<\infty\), and \(m\) is a predictor.

\paragraph{Predictors with \(E f(X)^2=\infty\).}
From \((a+b)^2\leq2a^2+2b^2\),
\[
 f(X)^2\leq2(f(X)-Y)^2+2Y^2,\qquad
 f(X)^2\leq2(f(X)-m(X))^2+2m(X)^2 .
\]
Because \(EY^2\) and \(E\,m(X)^2\) are finite, both \(R(f)=E(Y-f(X))^2\) and
\(E(f(X)-m(X))^2\) are infinite. The right side of the claimed identity is
\(EU^2+\infty=\infty\), so the identity holds.

\paragraph{Predictors with \(E f(X)^2<\infty\).}
Put \(V=m(X)-f(X)\); then \(EV^2<\infty\) and \(V\) is
\(\sigma(X)\)-measurable. Pointwise
\[
 (Y-f(X))^2=(U+V)^2=U^2+2UV+V^2,
\]
and \(|UV|\leq(U^2+V^2)/2\), so every term is integrable. By
\cref{thm:p15-conditional-basic}, \(C_X(m(X))=m(X)\) because \(m(X)\) is
\(\sigma(X)\)-measurable, so by linearity \(C_XU=C_XY-m(X)=0\) almost surely.
The variables \(U\) and \(VU\) are integrable and \(V\) is
\(\sigma(X)\)-measurable, so \cref{thm:p15-conditional-pullout} gives
\[
 C_X(VU)=V\,C_XU=0\quad\text{almost surely,}
\]
and the identity \(E(C_XW)=EW\) of \cref{thm:p15-conditional-basic} gives
\(E(UV)=0\). Therefore \(R(f)=EU^2+EV^2\), which is the claimed
identity.

\paragraph{Consequences.}
For every predictor, \(R(f)\geq EU^2=R(m)\), the last equality being the
identity for \(f=m\). Since \(m\) is itself a predictor,
\(R^{*}=R(m)=EU^2<\infty\). If \(R(f)=R^{*}\), then \(R(f)<\infty\), so
\(E f(X)^2<\infty\) by the second paragraph and
\(E(f(X)-m(X))^2=R(f)-R^{*}=0\), that is \(f(X)=m(X)\) almost surely.
Conversely, if \(f(X)=m(X)\) almost surely, the identity gives
\(R(f)=EU^2=R^{*}\). The excess risk is \(R(f)-R^{*}=E(f(X)-m(X))^2\),
including the value \(\infty\).
\end{proof}
\end{annexitem}

\begin{annexitem}{annex:p15-w02b-p15-064-classification-bayes}{thm:p15-w02b-p15-064-classification-bayes}{Risk as an integral of the conditional error probability}
\begin{proof}
\paragraph{A version with values in \([0,1]\).}
\Cref{thm:p15-conditional-factorization} gives a measurable real
\(\eta_0\) with \(\eta_0(X)\) a version of
\(P(Y=1\mid X)=C_X\mathbf1_{\{Y=1\}}\). By the order property of
\cref{thm:p15-conditional-basic}, \(0\leq\eta_0(X)\leq1\) almost surely, so
\(\eta=\min\{\max\{\eta_0,0\},1\}\) is measurable, has values in \([0,1]\),
and \(\eta(X)=\eta_0(X)\) almost surely; hence \(\eta(X)\) is also a version.
The theorem applies to any such choice.

\paragraph{The risk of a predictor.}
Let \(f\) be a predictor and \(A=\{f(X)=+1\}\in\sigma(X)\). Pointwise
\[
 \mathbf1_{\{f(X)\neq Y\}}
 =\mathbf1_A\mathbf1_{\{Y=-1\}}+\mathbf1_{A^c}\mathbf1_{\{Y=+1\}} .
\]
The defining property of the version \(\eta(X)\), applied on the events
\(A\) and \(A^c\) of \(\sigma(X)\), gives
\(E(\mathbf1_A\mathbf1_{\{Y=1\}})=E(\mathbf1_A\eta(X))\) and
\(E(\mathbf1_{A^c}\mathbf1_{\{Y=1\}})=E(\mathbf1_{A^c}\eta(X))\); since
\(\mathbf1_{\{Y=-1\}}=1-\mathbf1_{\{Y=1\}}\),
\(E(\mathbf1_A\mathbf1_{\{Y=-1\}})=E(\mathbf1_A(1-\eta(X)))\). Adding,
\[
 R(f)=E\bigl[\mathbf1_A(1-\eta(X))+\mathbf1_{A^c}\eta(X)\bigr]
 =E\,\eta(X)+E\bigl[\mathbf1_A(1-2\eta(X))\bigr].
 \tag{1}
\]

\paragraph{The excess-risk formula.}
\(f_{\mathrm{Bayes}}\) is measurable because \(\{\eta\geq1/2\}\) is a
measurable set; it is a predictor. Let \(B=\{f_{\mathrm{Bayes}}(X)=+1\}=
\{\eta(X)\geq1/2\}\). All risks here lie in \([0,1]\), so (1) for \(f\) and
for \(f_{\mathrm{Bayes}}\) may be subtracted:
\[
 R(f)-R(f_{\mathrm{Bayes}})=E\bigl[(\mathbf1_A-\mathbf1_B)(1-2\eta(X))\bigr].
\]
On \(A\setminus B\), \(\mathbf1_A-\mathbf1_B=1\) and \(\eta(X)<1/2\), so the
integrand is \(1-2\eta(X)=|2\eta(X)-1|\). On \(B\setminus A\),
\(\mathbf1_A-\mathbf1_B=-1\) and \(\eta(X)\geq1/2\), so the integrand is
\(2\eta(X)-1=|2\eta(X)-1|\). Elsewhere the integrand is zero. Since
\((A\setminus B)\cup(B\setminus A)=\{f(X)\neq f_{\mathrm{Bayes}}(X)\}\), the
formula follows, and its right side is nonnegative.

\paragraph{The minimum and the minimizers.}
By the formula, \(R(f)\geq R(f_{\mathrm{Bayes}})\) for every predictor, so
\(R^{*}=R(f_{\mathrm{Bayes}})\). By (1) with \(A=B\),
\[
 R(f_{\mathrm{Bayes}})=E\bigl[\mathbf1_B(1-\eta(X))+\mathbf1_{B^c}\eta(X)\bigr]
 =E\min\{\eta(X),1-\eta(X)\},
\]
because \(1-\eta(X)\leq\eta(X)\) on \(B\) and \(\eta(X)<1-\eta(X)\) on
\(B^c\). A predictor \(f\) has \(R(f)=R^{*}\) exactly when the nonnegative
variable \(|2\eta(X)-1|\mathbf1_{\{f(X)\neq f_{\mathrm{Bayes}}(X)\}}\) has
expectation zero, that is, vanishes almost surely; since \(|2\eta(X)-1|>0\)
exactly on \(\{\eta(X)\neq1/2\}\), this says that \(f(X)=f_{\mathrm{Bayes}}(X)\)
almost surely on that event.

\paragraph{Labels \(\{0,1\}\).}
Let \(\varphi(0)=-1\), \(\varphi(1)=+1\). For a record \((X,Y)\) with
\(Y\in\{0,1\}\) and a predictor \(g\) with values in \(\{0,1\}\),
\(\mathbf1_{\{g(X)\neq Y\}}=\mathbf1_{\{\varphi(g(X))\neq\varphi(Y)\}}\), so
the risk of \(g\) for \((X,Y)\) is the risk of \(\varphi\circ g\) for
\((X,\varphi(Y))\). Composition with \(\varphi\) is a bijection between the
two sets of predictors, and \(P(\varphi(Y)=1\mid X)=P(Y=1\mid X)\). Every
statement proved above therefore transfers, with \(0\) in place of \(-1\).
\end{proof}
\end{annexitem}

\begin{annexitem}{annex:p15-w02b-p15-064-decomposition}{thm:p15-w02b-p15-064-decomposition}{The approximation term is a constant of the problem}
\begin{proof}
The displayed identity is algebra in \((-\infty,\infty]\), the subtracted
quantities \(R_{\mathcal F}\) and \(R^{*}\) being finite. For \(f\in\mathcal F\),
\(R(f)\geq R_{\mathcal F}\) by definition of the infimum, and
\(R_{\mathcal F}\geq R^{*}\) because \(\mathcal F\subseteq\mathcal F_{\mathrm{all}}\)
and an infimum over a larger set is no larger. \(R_{\mathcal F}\) and
\(R^{*}\) are defined from the risks of fixed predictors, which depend only on
\(\ell\) and \(P\), so their difference involves no sample.

Item~1: for each sample, \(\widehat f_n\) is a member of \(\mathcal F\), so
\(R(\widehat f_n)\geq R_{\mathcal F}\); subtract \(R^{*}\). Item~2: add
\(R_{\mathcal F}-R^{*}\) to both sides of the assumed inequality and use the
identity. Item~3: if \(0<\varepsilon<R_{\mathcal F}-R^{*}\), item~1 gives
\(R(\widehat f_n)-R^{*}>\varepsilon\) for every sample, so the event has
probability one for every \(n\), and the probabilities do not tend to zero.
Item~4: if \(R_{\mathcal F}=R^{*}\), then \(R(\widehat f_n)-R^{*}\) and
\(R(\widehat f_n)-R_{\mathcal F}\) are the same random variable, so the two
families of events in \cref{def:p15-w02b-p15-064-consistency} coincide.
\end{proof}
\end{annexitem}

\begin{annexitem}{annex:p15-w02b-p15-064-finite-class-bayes}{thm:p15-w02b-p15-064-finite-class-bayes}{Combining the uniform event with the decomposition}
\begin{proof}
Each \(h\in\mathcal H\) is a predictor, and
\(z=(x,y)\mapsto\ell(x,y,h(x))\) is measurable with values in \([0,1]\), so
\cref{def:p15-conc-learning} holds and its risks are those of
\cref{def:p15-w02b-p15-064-risk}. Hence
\(R^{*}\leq\min_hR(h)\leq1<\infty\), and \(\min_hR(h)=R_{\mathcal H}\), a
minimum of finitely many numbers.

\paragraph{The two-sided bound.}
On the event of \cref{thm:p15-conc-finite-learning},
\(R(\widehat h)\leq\min_hR(h)+2b_n+\eta\); subtracting \(R^{*}\) gives the
upper bound. The lower bound holds for every sample by item~1 of
\cref{thm:p15-w02b-p15-064-decomposition}.

\paragraph{Measurability.}
The exact minimizer with the order tie rule of
\cref{def:p15-conc-learning} is determined by finitely many comparisons
between the measurable variables \(\widehat R_n(h_j)\), so each event
\(\{\widehat h=h_j\}\) is measurable and
\(R(\widehat h)=\sum_jR(h_j)\mathbf1_{\{\widehat h=h_j\}}\) is a random variable.

\paragraph{Consistency relative to \(\mathcal H\).}
Fix \(\varepsilon>0\) and \(0<\delta<1\). For
\(n\geq\lceil2\ln(2r/\delta)/\varepsilon^2\rceil\),
\cref{thm:p15-conc-finite-learning} gives
\(P\{R(\widehat h)-\min_hR(h)>\varepsilon\}\leq\delta\). Hence the
\(\limsup\) of these probabilities is at most \(\delta\) for every
\(\delta\in(0,1)\), and so it is zero. No property of \(P\) was used.

\paragraph{Bayes-consistency.}
If some \(h\in\mathcal H\) has \(R(h)=R^{*}\), then
\(R_{\mathcal H}=R^{*}\), and item~4 of
\cref{thm:p15-w02b-p15-064-decomposition} turns consistency relative to
\(\mathcal H\) into Bayes-consistency. If no member has risk \(R^{*}\), then,
the minimum over the finite class being attained by some member,
\(R_{\mathcal H}>R^{*}\), and item~3 shows that the procedure is not
Bayes-consistent.
\end{proof}
\end{annexitem}

\begin{annexitem}{annex:p15-w02b-p15-064-squared-example}{thm:p15-w02b-p15-064-squared-example}{Moments of the uniform law and the sample-mean minimizer}
\begin{proof}
\paragraph{The regression function and the minimum.}
\(EX^2=\int_0^1x^2\,dx=1/3\) and \(EX^4=\int_0^1x^4\,dx=1/5\). \(Y\) is
integrable and \(Y^2\leq2X^4+2\sigma^2\xi^2\), so \(EY^2<\infty\). The
variable \(\xi\) is integrable and independent of \(\sigma(X)\), so
\(C_X\xi=E\xi=0\) by \cref{thm:p15-conditional-basic}, and \(C_X(X^2)=X^2\)
because \(X^2\) is \(\sigma(X)\)-measurable. By linearity
\(C_XY=X^2\) almost surely, so \(m(x)=x^2\) is a regression function. By
\cref{thm:p15-w02b-p15-064-squared-bayes},
\(R^{*}=E(Y-X^2)^2=\sigma^2E\xi^2=\sigma^2\).

\paragraph{Constant predictors.}
For the constant predictor \(c\), the same theorem gives
\(R(c)=\sigma^2+E(X^2-c)^2\), and
\[
 E(X^2-c)^2=EX^4-2c\,EX^2+c^2=\tfrac15-\tfrac{2c}3+c^2
 =\bigl(c-\tfrac13\bigr)^2+\tfrac15-\tfrac19
 =\bigl(c-\tfrac13\bigr)^2+\tfrac4{45}.
\]
The infimum over \(c\) is attained only at \(c=1/3\), with value
\(\sigma^2+4/45\).

\paragraph{The empirical minimizer.}
With \(\overline Y_n=n^{-1}\sum_iY_i\) and
\(\sum_i(Y_i-\overline Y_n)=0\),
\[
 \widehat R_n(c)=\frac1n\sum_{i=1}^n(Y_i-c)^2
 =\frac1n\sum_{i=1}^n(Y_i-\overline Y_n)^2+(\overline Y_n-c)^2,
\]
whose unique minimizer is \(c=\overline Y_n\). Evaluating the risk formula at
this value, \(R(\overline Y_n)-R_{\mathcal F_c}=(\overline Y_n-1/3)^2\) and
\(R(\overline Y_n)-R^{*}=4/45+(\overline Y_n-1/3)^2\geq4/45\), for every
sample.

\paragraph{The expected in-class excess.}
\(EY=EX^2+\sigma E\xi=1/3\), and \(\operatorname{Var}Y=E(Y-1/3)^2=R(1/3)
=\sigma^2+4/45\). The \(Y_i\) are independent with the law of \(Y\), so, as
in the proof of \cref{thm:p15-w02b-p15-064-empirical-estimate} with
\(D_i=Y_i-1/3\),
\(E(\overline Y_n-1/3)^2=\operatorname{Var}Y/n=(\sigma^2+4/45)/n\).
\end{proof}
\end{annexitem}

\begin{annexitem}{annex:p15-w02b-p15-064-classification-example}{thm:p15-w02b-p15-064-classification-example}{The Bayes classifier of a two-point input}
\begin{proof}
The input takes the two values \(0,1\), each with probability \(1/2\). By
\cref{thm:p15-conditional-partitions}, a version of \(P(Y=1\mid X)\) is
\(\eta(X)\) with \(\eta(0)=1/4\) and \(\eta(1)=3/4\). By
\cref{thm:p15-w02b-p15-064-classification-bayes} in its \(\{0,1\}\) form,
the Bayes classifier predicts \(1\) exactly where \(\eta\geq1/2\), that is
at \(x=1\), and \(0\) at \(x=0\); it is the identity predictor, and
\[
 R^{*}=E\min\{\eta(X),1-\eta(X)\}=\tfrac12\cdot\tfrac14+\tfrac12\cdot\tfrac14=\tfrac14 .
\]
The constant predictor \(0\) errs exactly when \(Y=1\), with probability
\(\tfrac12(\tfrac14+\tfrac34)=\tfrac12\); the constant \(1\) errs exactly
when \(Y=0\), also with probability \(\tfrac12\). Hence
\(\min_{h\in\mathcal H}R(h)=1/2\), and any predictor selected from
\(\mathcal H\) has risk \(1/2\): excess \(0\) relative to \(\mathcal H\) and
excess \(1/4\) over \(R^{*}\). The four-predictor class of
\cref{thm:p15-conc-noisy-example} contains the identity predictor, so its
minimum risk is \(R^{*}\), and item~4 of
\cref{thm:p15-w02b-p15-064-decomposition} shows that excess relative to it
and excess over \(R^{*}\) are the same quantity.
\end{proof}
\end{annexitem}

\begin{annexitem}{annex:p15-w02b-p15-064-failures}{thm:p15-w02b-p15-064-failures}{A memorizing predictor and an output without second moment}
\begin{proof}
\paragraph{Measurability of the memorizing predictor.}
Write
\[
 s=((x_1,y_1),\ldots,(x_n,y_n))
\]
for a realized sample. Then
\[
 \widehat f_s(x)=\sum_{i=1}^ny_i\,\mathbf1_{\{x=x_i,\ x\neq x_j\ (j<i)\}}
 +\mathbf1_{\{x\notin\{x_1,\ldots,x_n\}\}},
\]
a finite combination of indicators of sets defined by equalities of
coordinates, which are Borel in \((x,s)\). So \(\widehat f\) is jointly
measurable, each \(\widehat f_s\) is a predictor, and \(R(\widehat f_n)\) is
measurable by Tonelli's theorem.

\paragraph{Empirical risk.}
For \(i\neq j\), \(X_i\) and \(X_j\) are independent and the uniform law has
no atoms, so by Tonelli's theorem
\(P(X_i=X_j)=\int P(X_i=x)\,P_{X_j}(dx)=0\). Off the finite union of these
null events the inputs are distinct, the smallest index \(k\) with
\(X_k=X_i\) is \(i\), and \(\widehat f_n(X_i)=Y_i\); every training loss is
zero. Hence \(\widehat R_n(\widehat f_n)=0\) almost surely.

\paragraph{Risk.}
Fix a realized sample \(s\) and a fresh record \((X,Y)\). The finite set
\(\{x_1,\ldots,x_n\}\) has uniform probability zero, so
\(\widehat f_s(X)=+1\) almost surely and
\(R(\widehat f_s)=P(Y\neq+1)=1/2\). Since \(Y\) is independent of \(X\),
\cref{thm:p15-conditional-basic} gives the constant version
\(\eta\equiv P(Y=1)=1/2\), and
\cref{thm:p15-w02b-p15-064-classification-bayes} gives
\(R^{*}=E\min\{1/2,1/2\}=1/2\). Thus \(R(\widehat f_n)=1/2=R^{*}\) for every
sample, and the two expectations are \(0\) and \(1/2\).

\paragraph{Infinite minimum risk.}
Independence of \(X\) and \(Y\) means that the joint law agrees with the
product \(P_X\otimes P_Y\) on measurable rectangles; rectangles form a
pi-system generating the product sigma-algebra, so the two probability
measures agree on it. For a predictor \(f\), Tonelli's theorem gives
\[
 R(f)=\int_{\mathcal X}\Bigl(\int_{\mathbb R}(y-f(x))^2\,P_Y(dy)\Bigr)P_X(dx).
\]
For every real \(c\), \(y^2\leq2(y-c)^2+2c^2\), so
\(\infty=EY^2\leq2E(Y-c)^2+2c^2\) and \(E(Y-c)^2=\infty\). The inner integral
is therefore infinite at every \(x\), and \(R(f)=\infty\). As \(f\) was
arbitrary, \(R^{*}=\infty\).
\end{proof}
\end{annexitem}
